\begin{python0}
from solutions import *; clear()
\end{python0}
\sectionthree{C-strings}

A \EMPHASIZE{C-string} is nothing more than an array of characters. The only thing different between a C-string and an array of integers (say) is that the \textbf{special escape character '\ 0'} (the 0 is the \textbf{number} 0, not the letter O) is used to mark the \textbf{end-of-string}. This is a special character just like \verb!'\n'! or \verb!'\t'! because \texttt{'\textbackslash 0'}, \verb!'\n'!, \verb!'\t'! are not ``visible'' characters.

Try to print \texttt{'\textbackslash 0'} and you won't see a thing:

\begin{consolethree}[escapeinside=||]
std::cout << '\0' << std::endl;
\end{consolethree}

(There are in fact lots of invisible characters.)

The point of the end-of-string marker \texttt{'\textbackslash 0'}, is to end some string processing. This is what I mean. Try this:

\begin{consolethree}[escapeinside=||]
char spam[] = {'a', 'b', 'c', '\0', 'd', 'e'};
std::cout << spam << std::endl;
\end{consolethree}

In this case the ``processing'' is the print statement:

\begin{consolethree}[escapeinside=||]
std::cout << spam << std::endl;
\end{consolethree}

The \texttt{'\textbackslash 0'} basically tells the print statement that printing should stop after printing the third character \texttt{'c'}. The output looks like this:

\begin{center}
\texttt{abc}
\end{center}

Note another thing ... \EMPHASIZE{You can actually print a character array!!!} Remember that you cannot do that for an array of integers, doubles, or booleans. You have to write a for-loop in those case. As a reminder try this:

\begin{consolethree}[escapeinside=||]
int x[] = {1, 2, 3, 4, 5};
std::cout << x << std::endl;
\end{consolethree}

Let's go back to the above example:

\begin{consolethree}[escapeinside=||]
char spam[] = {'a', 'b', 'c', '\0', 'd', 'e'};
std::cout << spam << std::endl;
\end{consolethree}

Note that it does NOT mean that the size of the array is 3; it is true that three characters are printed. There are actually 6 characters in the array:

\begin{center}
\texttt{'a', 'b', 'c', \texttt{'\textbackslash 0'}, 'd', 'e'}
\end{center}

So you need to distinguish the difference between the \EMPHASIZE{\texttt{size}} of the string \texttt{spam} and its \EMPHASIZE{\texttt{length}} which is defined to be the number of characters up to, but not including, the first \texttt{'\textbackslash 0'}.

Informally when we think of

\begin{consolethree}[escapeinside=||]
char spam[] = {'a', 'b', 'c', '\0', 'd', 'e'};
\end{consolethree}

\EMPHASIZE{as a string}, we say that it is made up of 'a', 'b', 'c'. But \EMPHASIZE{as a character array} is made up of six characters. Make sure you remember that.

You can think of the character \texttt{'\textbackslash 0'} in the array

\begin{consolethree}[escapeinside=||]
char spam[] = {'a', 'b', 'c', '\0', 'd', 'e'};
\end{consolethree}

as a sentinel value to mark the end of data for the string represented as an array. The character '\ 0' is also called the \EMPHASIZE{null character}. This is the reason why C-strings are also called \EMPHASIZE{null-terminated strings}.

\begin{ex}
What are the lengths and sizes of the following variables?
\begin{consolethree}[escapeinside=||]
char ham[] = {'3', '.', '\0', '1', '4'};
char eggs[] = {'\0', 'b', '\0', 'a'};
\end{consolethree}

You can verify your answers by counting the number of characters printed:
\begin{consolethree}[escapeinside=||]
std::cout << ham << std::endl;
std::cout << eggs << std::endl;
\end{consolethree}
\end{ex}

There is a shorthand for \{'a', 'b', 'c', '\ 0'\}. Try this:

\begin{consolethree}[escapeinside=||]
char spam[] = "abc"
std::cout << spam << std::endl;
\end{consolethree}

Note in this case that spam has size 4 (there are 4 characters in the array) and length 3 (there are 3 characters before the first \texttt{'\textbackslash 0'}); "abc" is

\begin{center}
'a', 'b', 'c', \texttt{'\textbackslash 0'}.
\end{center}

\begin{ex}
What is the length and size of the following variable?
\begin{consolethree}[escapeinside=||]
char spam[100] = "ab\0cd";
\end{consolethree}
\end{ex}

\begin{ex}
What is the length and size of the following variable?
\begin{consolethree}[escapeinside=||]
char spam[] = "ab\0cd";
\end{consolethree}
\end{ex}

\begin{ex}
What are the length and size of the following variable?
\begin{consolethree}[escapeinside=||]
char spam[100] = "";
\end{consolethree}
\end{ex}

Of course you can do this:

\begin{consolethree}[escapeinside=||]
std::cout << "abc\0de" << std::endl;
\end{consolethree}

You have already seen this ...

\begin{consolethree}[escapeinside=||]
std::cout << "hello world" << std::endl;
\end{consolethree}

during the first week of class. Now you know that \texttt{"hello world"} is an array of characters. It's just the character array made up of the following characters:

\begin{center}
\texttt{'h','e','l','l','o',' ','w','o','r','l','d','\textbackslash 0'}
\end{center}

(Don't forget the \texttt{'\textbackslash 0'}!!!!)

\begin{ex}
Can you do the following?
\begin{consolethree}[escapeinside=||]
std::cout << {'a', 'b', 'c'} << std::endl;
\end{consolethree}
\end{ex}

If you have a long string you can break it up like this:

\begin{consolethree}[escapeinside=||]
char foo[] = "abc"
             "def";
std::cout << foo << std::endl;
\end{consolethree}

Basically C++ will join it up for you so that the above code is the same as this:

\begin{consolethree}[escapeinside=||]
char foo[] = "abcdef";
std::cout << foo << std::endl;
\end{consolethree}

\begin{ex}
Can you do this:
\begin{consolethree}[escapeinside=||]
std::cout << "abc" "def" << std::endl;
\end{consolethree}

Can you do the same for integers?
\begin{consolethree}[escapeinside=||]
std::cout << 123 456 << std::endl;
\end{consolethree}
\end{ex}

\begin{ex}
You now know that although you cannot print an array of integers with \texttt{std::cout} (you need to write a for-loop), you can print an array of characters (it will only print up to
\texttt{'\ 0'}). What about input? Try this:
\begin{consolethree}[escapeinside=||]
char s[100];
std::cin >> s; // type in abcde
std::cout << s << std::endl;
\end{consolethree}

Not too shocking right?
\end{ex}

\begin{ex}
Run the above program again and enter this as input:

\begin{center}
\texttt{hello world}
\end{center}

What do you see?
\end{ex}

The problem with \texttt{std::cin} is that is that \texttt{std::cin} cuts up inputs at spaces, tabs, or newline (i.e. whitespaces). So for the above exercise, C++ will only give the string \texttt{"hello"} to variable s. If you do want spaces to go into your string variable, then you need to do something else for input. Check out a later section on the \texttt{getline()} function.


